% New main-task proof. Root owns A.4.54; no prior source is overwritten.
\subsection{Finite recognition of the actual involution fixed-locus loops}

For $9\leq n\leq12$, put $k=n-8$ and take the labelled torus
$T=\mathbb R^2/\mathbb Z^2$ with its nine points $T[3]$ and
$k-1$ additional, individually named points chosen from
$T[2]\setminus\{0\}$.  The distinguished label $0$ is their
common origin.  The actual map $z\mapsto-z$ induces a labelled
involution $\iota$: it fixes the $k$ named two-torsion points and
exchanges the points in each of four named pairs.
Write
\[
 \begin{gathered}
 X_n=\operatorname{Conf}_n(T)/T
       \xrightarrow{\sim}\operatorname{Conf}_{n-1}(T\setminus\{0\}),
 \qquad (z_i)\longmapsto(z_i-z_0)_{i\ne0},\\
 K_n=\pi_1(X_n)=\ker\bigl(P\Gamma_{1,n}\to
                                      \operatorname{SL}_2(\mathbb Z)\bigr),
 \qquad g_n=\sum_{j=2}^{n}j=\frac{n(n+1)}2-1.
 \end{gathered}                                                 \tag{ICR.1}
\]
The chosen torsion configuration is the basepoint $x_n$.  If $\sigma$
is the named permutation of the four pairs and the $k$ fixed labels,
the action of $\iota$ on this actual normalized space is
\[
                         (\iota z)_i=-z_{\sigma(i)}.
                                                               \tag{ICR.2}
\]
Let $X_n^{\iota,0}$ denote the connected component of its fixed
locus containing $x_n$.  The inclusion of this component defines
the based group homomorphism
\[
              a_n:\pi_1(X_n^{\iota,0},x_n)\longrightarrow K_n.
                                                               \tag{ICR.3}
\]

Start with a specified actual finite quotient
$q_0:\mathcal M_n=\Gamma_{1,[n]}\twoheadrightarrow Q_0$.
For the actual marked set $Z_n\subset T$, retain the relative-homology
representation
\[
 r_2:\mathcal M_n\longrightarrow
          \operatorname{GL}(H_1(T,Z_n;\mathbb F_2)),\qquad
 \dim_{\mathbb F_2}H_1(T,Z_n;\mathbb F_2)=n+1.
                                                               \tag{ICR.14}
\]
It is induced by the actual action on the topological pair $(T,Z_n)$.
Let $q=(q_0,\operatorname{lab},r_2)$ with target its actual image $Q$.
Put $d=|Q|\leq |Q_0|n!2^{(n+1)^2}$ and
$U_0=\ker(q|_{K_n})$.  Define the following explicit integers:
\[
 \begin{split}
 A_2(r)&=2^{\,r(r+1+2^r(r-1))},\qquad C_2(1,r)=2^r,\\
 a&=C_2(l-1,r)A_2(r),\\
 C_2(l,r)&=a\,C_2\bigl(l-1,1+a(r-1)\bigr)
                        A_2\bigl(1+a(r-1)\bigr)\quad(l\geq2),\\
 d_0&=d,\qquad
 m_i=d_i C_2\bigl(n-1,1+(n-1)d_i\bigr),\\
 d_{i+1}&=(m_i!)^{(m_i!)^{g_n}}\quad(0\leq i<2n-1),\\
 U_{i+1}&=\kappa_{m_i}(K_n),\qquad V=U_{2n-1}.
 \end{split}                                                    \tag{ICR.4}
\]
As before, $\kappa_m$ is the intersection of the kernels of all
homomorphisms to $\mathfrak S_m$, on the displayed ordinary group.

\begin{result}{A.4.54}{Proved}{An explicit binary refinement recognizes the ordinary involution centralizer in each configuration kernel}
The fixed locus of (ICR.2) has respectively $1,3,6,6$ connected
components for $n=9,10,11,12$.  Each component has real dimension
eight, and its ordered fixed-label assignment is its exact component
invariant.  The actual map (ICR.3) has image
\[
                        \operatorname{im}a_n=C_{K_n}(\iota).
                                                               \tag{ICR.5}
\]
For the explicit refinement (ICR.4), the following equality holds
under the natural projection of the two specified semidirect products:
\[
 \operatorname{im}\left(
 C_{(K_n/V)\rtimes\langle\iota\rangle}(\iota)\cap K_n/V
                              \longrightarrow K_n/U_0\right)
               =\operatorname{im}\bigl(C_{K_n}(\iota)\to K_n/U_0\bigr).
                                                               \tag{ICR.6}
\]
Moreover $[K_n:V]\leq d_{2n-1}$.  There is an actual finite quotient
$\eta_n:\mathcal M_n\to E_n$ with a projection to $Q$ such that
\[
 \begin{gathered}
 \ker\eta_n\cap K_n=V,\qquad \eta_n(K_n)=K_n/V,\\
 |E_n|\leq I_n\bigl(d\,d_{2n-1}^{g_n+1}\bigr)^{I_n},
                  \qquad I_n=144\,n!.
 \end{gathered}                                                 \tag{ICR.7}
\]
The identification in the first line is induced by the identity on
$K_n$.  Write $p:E_n\to Q$ for the stated projection.  The named
involution embeds, and the entire ambient centralizer has the exact
ordinary image
\[
                 p\bigl(C_{E_n}(\eta_n\iota)\bigr)
                            =q\bigl(C_{\mathcal M_n}(\iota)\bigr).
                                                               \tag{ICR.15}
\]
Projection onto $Q_0$ gives the same equality for that original
specified output.  Thus (ICR.15) includes the label and period
correction, as well as the kernel comparison (ICR.6).
\end{result}

\begin{proof}
\par\medskip\noindent\textbf{The fixed configurations and their components.}\quad
Every fixed labelled coordinate lies in $T[2]$.  Label $0$ is the
origin, and the other $k-1$ fixed labels must be assigned distinct
elements of the three-element set $T[2]\setminus\{0\}$.  Choose a
representative label in each of the four paired orbits, and denote
its coordinate by $u_j$.  Its mate has coordinate $-u_j$.
The exact remaining conditions are
\[
 u_j\notin T[2],\qquad u_i\ne u_j,\qquad u_i\ne-u_j
                                            \quad(i\ne j).
                                                               \tag{ICR.8}
\]
For each fixed-label assignment, forgetting the signs of these four
representatives gives a regular covering of degree $16$ onto
\[
       \operatorname{Conf}_4\bigl((T\setminus T[2])/\{\pm1\}\bigr).
                                                               \tag{ICR.9}
\]
The base surface in (ICR.9) is a sphere with its four branch points
removed.  It is connected, as is its ordered four-point configuration
space: move one point at a time along paths avoiding the finitely many
other points, perturbing the paths within that surface as necessary.
Moving the $j$th point once around one branch puncture, with the other
three points fixed, changes exactly the $j$th lift sign.  Thus the
monodromy contains the four independent generators of $C_2^4$ and
is transitive on the sixteen lifts.  The covering is connected.
This proves that the ordered fixed-label assignments are precisely
the connected components.  Their numbers are
$3!/(3-(k-1))!$, namely $1,3,6,6$.  The four free torus
representatives in (ICR.8) give real dimension eight.

These are literal fixed loci of the evaluation-space action.
In particular, forgetting the last added fixed label restricts on
them to forgetting that entry of the ordered assignment, with the
four coordinates $u_j$ retained.  The fibre counts on assignments
are consequently three, two and one in the successive additions.
No component is identified by changing one of its fixed labels.

\par\medskip\noindent\textbf{The ordinary centralizer as a fixed-loop image.}\quad
The normalized evaluation fibration gives (ICR.1), by the same
translation-factor map and inverse as (NCF.1).  The successive
forgettings of points have free kernels of ranks $2,3,\ldots,n$.
Their actual fibres are punctured tori, and their homotopy exact
sequences prove that $X_n$ is aspherical.  Its universal cover
$\widetilde X_n$ is therefore contractible.  The based lift of
(ICR.2) is the unique lift fixing the chosen point $\widetilde x_n$.
Smith theory for the smooth involution makes
$\widetilde X_n^{\iota}$ connected and mod-two acyclic
\cite[Theorem~3.3.1 and Proposition~3.3.2, p.~274]{ET_Edmonds2012}.
Its local fixed-set charts are smooth submanifold charts, so this
connected fixed set is path connected.

A loop in $X_n^{\iota,0}$ lifts to a path starting at
$\widetilde x_n$.  Applying $\iota$ gives the same lifted path,
by uniqueness of path lifting.  Its endpoint is a fixed point of
the form $v\widetilde x_n$, with $v\in K_n$.  The two lifts
$\iota$ and $v^{-1}\iota v$ then fix the same chosen point and
project to the same action; uniqueness gives $v\iota=\iota v$.
This proves $\operatorname{im}a_n\subset C_{K_n}(\iota)$.
Conversely a commuting $v$ carries $\widetilde x_n$ to another
fixed point.  Join them by a path in the connected fixed set and
project it.  The resulting loop lies in $X_n^{\iota,0}$ and has
the specified deck displacement $v$.  This proves (ICR.5).
It uses the image of (ICR.3); injectivity of that homomorphism is
not required.

\par\medskip\noindent\textbf{The binary cohomology refinement.}\quad
We give the prime-two ingredient explicitly.  For a free group $F$
of rank $r\geq2$, put $N=F'F^2$ and
$Q_F=F/(N'N^2)$.  Its order is
$2^{r+1+2^r(r-1)}$, from the actual $2^r$-sheeted rose cover.
The center of $Q_F$ projects trivially to $F/N=C_2^r$.
Indeed a nonzero translation $v\in C_2^r$ lies outside the line
spanned by at least one standard basis vector $e_j$.
The lift of the squared $j$th generator gives the nonzero
two-edge cycle supported on $\{0,e_j\}$.  Translation by $v$
moves it to the distinct parallel line.  Graph homology embeds in
edge chains, so these cycles are different in $H_1(N,\mathbb F_2)$.
Thus the action of $F/N$ on that homology is faithful, which proves
the asserted projection of the center.

For any actual extension with this normal free kernel, the kernel
of its conjugation map to $\operatorname{Aut}(Q_F)$ intersects $F$
inside $F'F^2$ and has index at most $|Q_F|^r=A_2(r)$.
In the induced Lyndon--Hochschild--Serre comparison, the $q=1$
coefficient map is zero on this intersection, while the positive
$q=0$ map factors through the chosen cohomology-killing subgroup of
the base.  There are at most two filtration pieces in each positive
total degree.  Performing this construction twice therefore kills
the actual positive cohomology restriction.  This is the same
two-row calculation (ECR.10), now with the displayed binary quotient
and coefficients $\mathbb F_2$.  Schreier's formula gives the second
rank bound $1+a(r-1)$ and hence exactly the recursion $C_2(l,r)$ in
(ICR.4), starting with the free case $[F:F'F^2]=2^r$.

An index-$d_i$ subgroup of $K_n$ inherits the actual length-$n-1$
free extension tower, with each free rank at most
$1+(n-1)d_i$.  The coset action on the constructed subgroup then
proves $U_{i+1}\subset U_i$ and that
\[
             H^{>0}(U_i,\mathbb F_2)\longrightarrow
             H^{>0}(U_{i+1},\mathbb F_2)\quad\hbox{is zero}.
                                                               \tag{ICR.10}
\]
The $g_n$ chosen source generators bound the index of
$\kappa_{m_i}(K_n)$ by the factorial expression in (ICR.4).
This proves every asserted index bound and kernel containment.
Each $U_i$ is invariant under $\iota$; all $U_i$ after $U_0$ are
characteristic in $K_n$.

\par\medskip\noindent\textbf{Fixed-component recognition at the finite quotient.}\quad
Use the periodic hyperplanes
$6(z_i^{(j)}-z_l^{(j)})\in\mathbb Z$, for $j=1,2$, including
the anchored zero coordinate.  They are invariant under (ICR.2),
have a finite cell quotient and contain every removed diagonal as
a subcomplex.  All chosen six-division basepoints are vertices.
After barycentric subdivision, deleting the forbidden diagonal
vertices gives the actual equivariant deformation retraction of
the diagonal complement used in (ECR.13).  Its finite complex has
dimension at most $2n-2$.  Lift it to the covers specified by the
$U_i$, and call the resulting finite complexes $Y_i$, with chosen
basepoint vertices $y_i$ fixed by the canonical lift of $\iota$.
Because the spaces are connected, (ICR.10) is zero in every degree
for cohomology relative to these basepoints.

The complete cyclic resolution for $C_2$ has alternating
differentials $1-\iota$ and $1+\iota$; both are zero on a trivial
$\mathbb F_2$-module.  The ordinary-cohomology-first Tate spectral
sequence for $C^*(Y_i,y_i;\mathbb F_2)$ has only $2n-1$ vertical
degrees.  Each map in (ICR.10) is zero on its $E_2$ page, so the
composite of $2n-1$ maps vanishes on the Tate abutment, by its
finite filtration in every total degree.  The relative cells
outside $Y_i^\iota$ occur in free $C_2$-orbits; their bounded Tate
complex is acyclic.  Restriction therefore identifies the abutment
with that on the fixed subcomplex.  On that fixed subcomplex the
two horizontal differentials are zero.  We obtain the actual zero
map
\[
 H^0(Y_0^\iota,\{y_0\};\mathbb F_2)
      \longrightarrow
 H^0(Y_{2n-1}^\iota,\{y_{2n-1}\};\mathbb F_2).
                                                               \tag{ICR.11}
\]
The characteristic functions of the components other than the one
containing $y_0$ show that every fixed point upstairs projects into
that one component downstairs.  This argument retains connected
components; it does not require a discrete fixed set.

Take $vV\in K_n/V$ centralizing the named $\iota$.  The fibre point
$vV\,y_{2n-1}$ is fixed.  Its lower fibre point $vU_0\,y_0$ therefore
lies in the component of $y_0$ in $Y_0^\iota$.  Join the two points
there by a path and project this path into $X_n^\iota$.
It is a based fixed-locus loop whose deck displacement in $K_n/U_0$
is $vU_0$.  Formula (ICR.5) therefore realizes this residue by an
ordinary centralizer element.  The reverse containment in (ICR.6)
is given by the images of the actual ordinary commuting elements.
This proves exactly the displayed equality.

\par\medskip\noindent\textbf{An actual finite ambient quotient.}\quad
Every marked point of the chosen base configuration lies in $T[6]$.
Consequently the literal linear maps with period in $\Gamma(6)$
fix all its labels, and give the homomorphic splitting
\[
 \mathcal U_n=P\Gamma_{1,n}\cap\operatorname{per}^{-1}\Gamma(6)
             =K_n\rtimes s\Gamma(6),\qquad
 [\mathcal M_n:\mathcal U_n]=144\,n!.
                                                               \tag{ICR.12}
\]
Here $144=|\operatorname{SL}_2(\mathbb Z/6)|=6\cdot24$ by the
literal Chinese remainder map; pure forgetting is onto the torus
group and disc braids give every label permutation with identity
compact period, proving the index.  The subgroup $V$ is
characteristic in the normal group $K_n$.  Define
\[
 \mathcal U_n\longrightarrow
 Q\times\bigl((K_n/V)\rtimes\operatorname{Aut}(K_n/V)\bigr),
 \qquad v s(A)\longmapsto
       \bigl(q(vs(A)),(vV,\operatorname{Ad}_{s(A)}|_{K_n/V})\bigr).
                                                               \tag{ICR.13}
\]
This is an actual homomorphism whose kernel intersects $K_n$
precisely in $V$.  Its image has order at most
$d[K_n:V]^{g_n+1}$, since images of the $g_n$ named generators
determine an automorphism of that finite group.
Induce (ICR.13) through the $I_n$ actual cosets of the normal
subgroup $\mathcal U_n$.  Every conjugate kernel still intersects
$K_n$ in $V$, giving the first line of (ICR.7).  Its permutation
image has order $I_n$, giving the second line.  The kernel of
induction is contained in $\ker q$, through the identity coset,
so it also supplies the required literal projection to $Q$.
The retained permutation of the four pairs detects the nonidentity
involution, so $\eta_n\iota$ has order two.  In particular its
subgroup with $\eta_n(K_n)$ is the specified semidirect product
in (ICR.6).

\par\medskip\noindent\textbf{The exact ordinary label and period correction.}\quad
Take $v\in C_{E_n}(\eta_n\iota)$ and choose an ordinary source
representative $g\in\mathcal M_n$.  Let $A\in\operatorname{SL}_2(\mathbb Z)$
be its actual compact period and let $\tau$ be its label permutation.
Since the quotient retains labels, $\tau$ commutes with the
involution permutation $\sigma$.  It therefore permutes the fixed
labels and the four paired orbits with their possible flips.
We now check the additional compatibility with $A$ using exactly
the retained factor (ICR.14).

Choose arcs $\gamma_j$ from the origin to each other marked point
$p_j$, and use their relative classes together with the two compact
cycles as a basis of $H_1(T,Z_n;\mathbb F_2)$.
The quotient relative classes correspond to
$e_j-e_0$ in $\widetilde H_0(Z_n;\mathbb F_2)$.
Write $P,S$ for the actual matrices of $\sigma,\tau$ on this
augmentation module.  The matrices of the involution and of $g$
then have the precise block forms
\[
 r_2(\iota)=\begin{pmatrix}I_2&t\\0&P\end{pmatrix},\qquad
 r_2(g)=\begin{pmatrix}\overline A&B\\0&S\end{pmatrix}.
                                                               \tag{ICR.16}
\]
Here the compact action of $\iota$ is $-I$, equal to $I$ modulo two.
For a fixed label $j$, the corresponding column is
\[
                           t_j=-2p_j\pmod2.
                                                               \tag{ICR.17}
\]
To verify (ICR.17), lift the arc to $\mathbb R^2$, with endpoint
equal to a chosen representative of $p_j$ plus an integral vector.
The difference of its image under $z\mapsto-z$ and the original
arc has compact displacement $-2p_j$ plus an even integral vector.
Thus its reduction is exactly (ICR.17), independently of the arc.
Set $t_0=0$ for the distinguished origin.

The commutation of the two matrices in (ICR.16) is the actual equation
\[
                        \overline A t+B P=B+t S.
                                                               \tag{ICR.18}
\]
For a fixed label $j$, $P(e_j-e_0)=e_j-e_0$, whereas
$S(e_j-e_0)=e_{\tau j}-e_{\tau0}$.  Applying (ICR.17)--(ICR.18)
to that column cancels $B$ and gives
\[
            A p_j+p_{\tau0}=p_{\tau j}\quad\hbox{in }T[2].
                                                               \tag{ICR.19}
\]
Thus the actual affine torus map
$\ell(z)=Az+p_{\tau0}$ sends every fixed labelled point to its
required image.  It commutes with $z\mapsto-z$ since
$2p_{\tau0}=0$, and has the prescribed integral period $A$.

The four images under $\ell$ of the original paired orbits are
four distinct points of the sphere with the four branch points
removed.  An isotopy of that sphere fixing the branch points moves
them to the required four original pair orbits in the order
prescribed by $\tau$.  The independent monodromies in (ICR.9)
allow the four lift signs to be prescribed as well.  Lift this
isotopy from the identity to the double covering torus.  Its endpoint
$\psi$ commutes with $\iota$, fixes all four two-torsion points,
has compact period identity, and gives the required four signs.
Consequently
\[
             \ell_0=\psi\ell\in C_{\mathcal M_n}(\iota),
 \qquad \operatorname{per}(\ell_0)=A,
 \qquad \operatorname{lab}(\ell_0)=\tau.
                                                               \tag{ICR.20}
\]
This construction is an ordinary isotopy and its specified lift;
it uses no finite or completed centralizer equality.

Now $g\ell_0^{-1}\in K_n$, and its image under $\eta_n$ commutes
with the named involution.  Apply (ICR.6) to this exact kernel
residue.  There is $c\in C_{K_n}(\iota)$ with
$q(g\ell_0^{-1})=q(c)$.  Therefore
$p(v)=q(g)=q(c\ell_0)$ is the image of an ordinary centralizer
element.  Conversely every ordinary centralizer element maps to
$C_{E_n}(\eta_n\iota)$.  This proves (ICR.15), and completes
all comparisons in the statement.
\end{proof}

The executed certificate
\path{check_involution_configuration_labels.py} constructs the
actual involution on the six-division coordinates and its relative
homology matrix, then tests every compact matrix modulo two and
every commuting label permutation.  Its complete counts are
\[
 \begin{array}{c|r|r|r}
 n&\text{fixed components}&\text{period--label cases}
                       &\text{compatible cases}\\\hline
 9&1&2304&2304\\
 10&3&4608&1536\\
 11&6&13824&2304\\
 12&6&55296&9216.
 \end{array}                                                   \tag{ICR.21}
\]
For each compatible case it constructs the actual missing block
$B$ in (ICR.16) and checks commutation of the full matrices.
For each other case it retains a nonzero fixed-column obstruction
to (ICR.18).  The $76032$ cases include all permutations of the four pairs
and their independent flips, as well as every motion of the
distinguished origin among the fixed labels.  The same program
constructs the $144$-element period image modulo six and its
Chinese remainder projections.  Its JSON, numerical archive and
output retain the component assignments, matrices, obstruction
columns and exact counts used above.
